\section{Color bridge: axis-triplet $U(3)$ structure and kinematic confinement}
\label{sec:color}

\subsection{Axis-aligned triplet and prestress-driven degeneracy}
\label{subsec:axis-triplet}

On the 6-neighbor axial-only cubic lattice it is natural to organize the per-wavepacket complex
envelope as a three-component triplet $\Psi=(\psi_x,\psi_y,\psi_z)^\top\in\C^3$ aligned with the
lattice axes.
The prestress parameter $\alpha$ controls the eigenvalue spread of the 3-channel Christoffel matrix
$D(k)$ without introducing off-diagonal coupling between axis components.

\paragraph{Non-degenerate limit ($\alpha\to 1$, no prestress).}
The lattice energy is purely Hookean longitudinal; $D(k)$ has maximally split eigenvalues.
The three axis channels are dynamically decoupled, each supporting an independent rank-1 Berry
connection $a_\mu^{(x)}, a_\mu^{(y)}, a_\mu^{(z)}$.
The gauge group is closest to $U(1)^3$.

\paragraph{Degenerate limit ($\alpha\to 0$, maximum prestress).}
The geometric-stiffness term $(1-\alpha)|\Delta u|^2$ dominates and contributes equal stiffness to
all components on every bond.
$D(k)\propto I$ becomes proportional to the identity at every $k$; all three lateral channels are
fully degenerate.
The transported subspace is genuinely 3-dimensional and the adiabatic-transport object is a
Wilczek--Zee connection $\mathcal{A}_\mu\in\mathfrak{u}(3)$ (\Cref{eq:wz-connection}).
At the project's default operating point $\alpha=0.2$ the lattice sits close to this degenerate
limit.

\paragraph{Role of microscopic anisotropy.}
The $U(3)$ gauge group on a 3-dimensional complex envelope is a generic representation-theoretic
fact about $\C^3$.
What the cubic anisotropy provides is the \emph{physically meaningful decomposition}: a preferred
3-dimensional subspace (the lateral triplet, distinct from $X^4$), a preferred axis-aligned basis
within it (which makes the $SU(3)$ traceless generators carry operational meaning as ``color'' labels
on solitons), and a natural identification of the $U(1)$ trace with the electromagnetism-like sector.
Without cubic anisotropy, no preferred basis exists and the $SU(3)$ generators would be arbitrary
linear combinations with no operational handle.
That anisotropy can carry genuine non-Abelian gauge structure in real space is established by
precedent: \citet{Chen2019NonAbelianOptics} demonstrate a synthetic $SU(2)$ gauge field for optical
waves from the off-diagonal permittivity of an anisotropic medium.

\subsection{\texorpdfstring{$U(3) \simeq U(1)\times SU(3)$ decomposition}{U(3) = U(1) x SU(3) decomposition}}
\label{subsec:u3-decomp}

Decomposing $\mathcal{A}_\mu\in\mathfrak{u}(3)$ into trace and traceless parts:
\begin{equation}
  \mathcal{A}_\mu
  \;=\;
  \underbrace{\tfrac{1}{3}\mathrm{tr}(\mathcal{A}_\mu)\,I_3}_{U(1)\text{ part}}
  \;+\;
  \underbrace{\sum_{a=1}^{8}A^a_\mu\,T^a}_{SU(3)\text{ part}},
  \label{eq:u3-decomp}
\end{equation}
where $T^a$ are the eight traceless $\mathfrak{su}(3)$ generators.
The $U(1)$ part behaves electromagnetism-like; the traceless part carries non-Abelian self-couplings
through the commutator term in $\mathcal{F}_{\mu\nu}$.

\paragraph{EM as the symmetric (trace) direction.}
The $U(1)$/trace component of $\Psi$ is the projection onto $(1,1,1)/\sqrt{3}$:
\begin{equation}
  \Psi_{\mathrm{EM}} \;=\;
  \frac{1}{\sqrt{3}}(\psi_x+\psi_y+\psi_z)\,\frac{(1,1,1)^\top}{\sqrt{3}}.
  \label{eq:em-trace}
\end{equation}
A pure-EM excitation displaces all three lateral components in phase with equal weight; the traceless
complement ($SU(3)$) encodes differential phase and amplitude between components.
A single-component displacement (e.g.\ $\psi_x\neq 0$, $\psi_y=\psi_z=0$) projects as $\tfrac13$
into the $U(1)$ trace and $\tfrac23$ into the $SU(3)$-traceless sector; it is not a pure-EM
excitation.

\paragraph{Charge-triality locking.}
Because $U(3)=\bigl(U(1)\times SU(3)\bigr)/\mathbb{Z}_3$ --- the center $\mathbb{Z}_3\subset SU(3)$
identified with the cube-roots of unity in $U(1)$ --- single-valuedness forces the trace charge $q$
and the $SU(3)$ triality $t\in\{0,1,2\}$ to satisfy $q\equiv -t\pmod{3}$, the standard relation for
$U(N)=\bigl(U(1)\times SU(N)\bigr)/\mathbb{Z}_N$ \citep{GeorgiGlashow1974,BaezHuerta2010GUTAlgebra}.
A color triplet therefore carries fractional ($\tfrac13$-integer) trace charge while a color singlet
carries integer charge --- charge fractionality is locked to color representation without reference to
dynamics.

\paragraph{Energetic sector coupling.}
The central-force anharmonicity (quartic coefficient $\propto k_s\alpha/a$, \Cref{eq:W4} of Paper~I)
couples the two sectors at quartic order through a vertex $\propto\alpha|\Psi_{\mathrm{tr}}|^2
|\Psi_\perp|^2$, which vanishes identically as $\alpha\to 0$.
The prestress parameter that splits the two sectors' spectral scales is therefore also the parameter
that couples them energetically.

\subsection{Topological structure of the two sectors}
\label{subsec:topology}

The $U(1)$ and $SU(3)$ factors carry sharply different topology
\citep{Kibble1976Topology,VilenkinShellard1994,MantonSutcliffe2004}:
\begin{itemize}
  \item $\pi_1(U(1))=\mathbb{Z}$ admits line vortex defects (codimension-2 in 3D), with phase
  winding by $2\pi$ around the core \citep{NielsenOlesen1973}.
  \item $\pi_1(SU(3))=0$ forbids vortices in the $SU(3)$ sector.
  \item $\pi_3(SU(3))=\mathbb{Z}$ admits Skyrmion/instanton-type textures (codimension-3 worldlines
  in 4D) \citep{Skyrme1962,AtiyahManton1989}.
\end{itemize}
The two sectors support topologically \emph{different} classes of localized objects: codimension-2
vortex worldsheets in the EM/$U(1)$ sector versus codimension-3 texture worldlines in the
color/$SU(3)$ sector.

However, because the full vacuum manifold is $U(3)$ --- not just $U(1)$ --- the $U(1)$ vortex is a
\emph{semilocal} vortex \citep{VachaspatiAchucarro1991Semilocal,AchucarroVachaspati2000}: $\pi_1$ of
the full $U(3)$ group is $\mathbb{Z}$ (the $SU(3)$ factor is simply connected), so the vortex is
topologically protected within the $U(1)$ factor but not against perturbations that mix in the $SU(3)$
sector.
Stability is therefore dynamical rather than purely topological, and may depend on the amplitude regime
and the ratio of core energy to orientational moduli
\citep{AuzziBolognesiEvslinKonishiYung2003,Eto2021OrientationalModuli}.

\subsection{Kinematic confinement of color}
\label{subsec:color-confinement}

A free linear triplet wavepacket propagating on the cubic lattice must satisfy two requirements:
\begin{description}
  \item[\normalfont Coherence.] The three components stay phase-aligned over long range.
  This requires the three branch frequencies to be degenerate: $g(\hat k)=0$, where
  $g(\hat k)=\sqrt{\sum_a(h_a-\bar h)^2}/\bar h$ measures the spread of the lateral branch
  frequencies in direction $\hat k$.
  On the 6-neighbor cubic lattice, $g(\hat k)=0$ holds \emph{only} along the body diagonal
  $\hat k\parallel[111]$.

  \item[\normalfont Color-activity.] Operationally meaningful $SU(3)$-traceless holonomy ---
  a finite fibre-internal rotation distinguishing the three axis components --- requires the branches
  to be \emph{non-degenerate}: $g(\hat k)\neq 0$, i.e.\ propagation \emph{off} $[111]$.
  Exactly on $[111]$ the degenerate triplet has no preferred internal basis and the color holonomy is
  undefined.
\end{description}

No linear propagation direction satisfies both simultaneously.
Off $[111]$ the unequal branch speeds dephase the triplet over long range; on $[111]$ there is no
resolvable color.
Consequently:
\begin{quote}
  \emph{Color charge cannot be carried to infinity by a free linear excitation.
  It has support only in a nonlinear, phase-locked bound state, where the three components are held
  together by the geometric coupling of Paper~I rather than by spectral degeneracy.}
\end{quote}
We call this \emph{kinematic confinement of color}.
It is a statement about asymptotic states (no colored free particle), not a dynamical derivation of a
linear inter-quark potential, string tension, or area law.

\paragraph{Falsifiable handle.}
The off-$[111]$ fibre-holonomy ratio scales with prestress as
\begin{equation}
  \frac{R(\alpha_2)}{R(\alpha_1)}
  \;=\;
  \frac{(3-2\alpha_2)/\alpha_2}{(3-2\alpha_1)/\alpha_1},
  \label{eq:holonomy-ratio}
\end{equation}
giving $0.3077$ for $(\alpha_1,\alpha_2)=(0.2,0.5)$.
A deviation above $\sim 10\%$ would falsify the spectral-susceptibility factorization underlying it.

\subsection{Numerical targets}
\label{subsec:numerical-targets}

The theory contains two distinct particle sectors with different topology and different numerical
questions.

\paragraph{Target 1: $U(1)$ carrier-phase vortex (electron-like sector).}
The $U(1)$ trace sector supports line vortices.
A preliminary rotating-frame-periodic JFNK relaxation of a pure-trace $Y_1^1$ vortex seed on a
$112^3$ lattice did not converge: the core drifted $\approx19$ lattice units off the symmetry axis,
and the measured winding changed sign relative to the seed.
This is the expected signature of semilocal orientational moduli with a dynamical $SU(3)$ sector.
A converged eigenstate has not been obtained; branch-selection machinery for stabilizing the semilocal
vortex remains future work.

\paragraph{Target 2: $\pi_3(SU(3))$ texture (baryon-like sector).}
The cubic substrate's axis-aligned triplet and $\pi_3(SU(3))=\mathbb{Z}$ admit topologically stable
Skyrmion-type textures.
The canonical ansatz is the \emph{hedgehog}
\begin{equation}
  \xi^i(\mathbf{x}) \;=\; f(r)\,\hat x^i,
  \label{eq:hedgehog}
\end{equation}
with $J=0$, $L=1$: internal index locked to spatial direction, invariant only under the combined
diagonal $SO(3)$.
The entire nontrivial structure lives in the $SU(3)$ traceless sector; the $U(1)$ trace averages to
zero on the sphere, giving a trace-neutral far field (neutron-like).
Topological stabilization ($B=1$ winding) is achieved by extending the hedgehog into a Skyrme-twisted
configuration.
A nonzero $L=0$ scalar admixture shifts the $U(1)$ far field to proton-like.

\subsection{Color bridge open problems}
\label{subsec:color-open}

\begin{enumerate}
  \item \textbf{Verify holonomy ratio \eqref{eq:holonomy-ratio} numerically.}
  Compute the off-$[111]$ fibre-holonomy at $\alpha=0.2$ and $\alpha=0.5$ and check the predicted
  ratio $0.3077$.

  \item \textbf{Converged $U(1)$ vortex eigenstate.}
  Develop branch-selection machinery to stabilize the semilocal vortex against orientational drift
  and obtain a converged periodic-BC eigenstate.

  \item \textbf{Baryon/hedgehog stability test.}
  Construct the hedgehog ansatz on the cubic lattice, evolve under periodic-BC dynamics, and
  determine whether a stable, localized $\pi_3$-winding mode exists.

  \item \textbf{Dynamical co-localization of $U(1)$ and $SU(3)$ sectors.}
  Determine whether the quartic coupling $\propto\alpha|\Psi_{\mathrm{tr}}|^2|\Psi_\perp|^2$
  is sufficient to dynamically bind the charge to the color core in a baryon-like state.
\end{enumerate}